\section{原创创新点拆解}

本研究的创新点并非凭空设想的技术堆砌，而是从《海洋大数据技术原理与应用》课程中“海洋大数据的计算引擎”这一核心议题出发，经历了一次完整的问题提出、实证证伪与问题锐化的演进过程。课上讲授的“海洋大数据的计算引擎”使作者联想到：自主水下航行器端侧的算力与能量十分有限，计算、测量与运动之间必须权衡。由此作者提出了第一个朴素问题：当航行器带着一张几天前的旧卫星图下水作业时，它究竟应该把时间花在机载继续推演计算上，还是花在停下来重新观测上？围绕这一问题，我们在东海赤潮巡测场景下开展了前期研究（第3节），结果却发现：在旧图上成百上千倍地增加推演计算只会让规划器在自己的想象模型里打分越来越高，真实收益早已饱和；而由于叶绿素场连续平滑，计算与观测的时间分配空间总共只有约0.02的对数浓度差距，远小于下水位置与地图新鲜度的影响。针对前期方法增益微弱、创新性不足的缺陷，作者对研究问题进行了关键的批判与锐化：既然多算只是把旧图上的答案算得更精致，那么当客观上存在好几个与现有旧图同样相容、却要求采取截然不同行动的海洋世界时，再多的计算也无法将它们区分开来；真正重要的问题应当放到强海流绕障这种“左右两条路、一张粗海流图、选错一侧就多走十几个小时”的刀刃上，直接去问：针对当前粗图推荐的航路，我们能不能在历史大数据界定的合理流场范围内，专门搜出一个让另一侧反而更好、且让当前选择吃大亏的合理反例？这一追问直接催生了本文的合理反例检验框架。表\ref{tab:innovation}从四个维度拆解了本文的核心创新点、所替代的已有做法、实证证据及适用局限。

\begin{table}[htbp]
\centering
\caption{本文核心原创创新点、已有做法对比、实验证据与局限拆解。}
\label{tab:innovation}
\small
\begin{tabular}{@{}p{0.15\linewidth}p{0.21\linewidth}p{0.23\linewidth}p{0.19\linewidth}p{0.14\linewidth}@{}}
\toprule
创新点 & 已有做法及其不足 & 本文做法 & 证据（图表） & 局限 \\
\midrule
1. 分离数值歧义与环境歧义的决策判据 & 确定性规划在粗图上盲目加密网格（误以为多算能消除一切模糊）；传统元推理只管何时停止搜索，不区分误差源于求解器还是环境模型失真。 & 提出数值歧义（小于粗精网格容限 $\tau_{\text{num}}$，靠按需精算消除）与环境歧义（合理集合内存在翻转决策的反例，只能靠新信息或保守应对）的二分准则。 & 图\ref{fig:overview}、图\ref{fig:curse}、图\ref{fig:ambiguity}(a)：64倍精算仅将选错率从21.7\%降至19.7\%，当天0.5°图则降至10.4\%。 & 仅在左右绕障的二元拓扑决策上验证，尚未扩展至多分支同伦路径。 \\
\addlinespace
2. 基于包络定理梯度的合理反例主动搜索 & 集合规划随机抽取数十个流场算均值或翻转比例；翻转比例只反映多少样本翻转了决定，不反映翻转后损失多大。 & 利用包络定理沿当前最快航线累加流速导数，一次梯度仅需2次最短路求解；先沿初始梯度跳至球面再投影细化，直接求最坏后悔。 & 图\ref{fig:main}、图\ref{fig:risk}(a)：最坏后悔前20\%任务捕获67\%损失，显著优于32成员集合翻转率（49\%）与地图差值（34\%）。 & 流场过大扰动下最短路拓扑可能跳变，使目标函数非凸、局部梯度近似产生折耗。 \\
\addlinespace
3. 海洋大数据驱动的合理集合与半径标定 & 鲁棒优化与信息间隙方法常假设空间均匀白噪声并主观设定不确定性半径，脱离海洋物理结构，易过度保守或漏报黑潮切变。 & 用三年历史再分析误差标定非均匀幅度图 $\sigma(\boldsymbol{x})$ 与76千米相关长度的无散流函数模态；以覆盖90\%以上真实两侧时间差标定最小半径 $r=1.5$。 & 图\ref{fig:data}、图\ref{fig:ambiguity}(b)、表\ref{tab:scope}：$r=1.5$在测试年份覆盖94\%真实差值；消融替换为历史相似日差值后保守增益消失（0.99 vs 0.95小时）。 & 依赖多年历史再分析数据标定平稳统计量，对台风过境等非平稳极端天气适应受限。 \\
\addlinespace
4. 权衡最坏后悔与观测代价的门控响应 & 自适应采样面向全场降方差去测流，或不论代价一律保守/一律上浮，忽视了原位测流看不见远方切变及上浮通信本身耗时。 & 以最坏后悔超阈值 $\theta=8$小时作为门控：超标且上浮代价低（$<1$小时）时上浮取新图，代价高或仅有原位测流时改选最坏后悔更小一侧。 & 图\ref{fig:main}、图\ref{fig:risk}(b)、表\ref{tab:scope}：门控保守将平均损失从0.98降至0.70小时（降28\%），并给出约1小时的上浮代价切换点。 & 不观测时门控保守（0.70小时）比一律保守（0.72小时）的额外边际增益较小，主收益源于最坏后悔准则本身。 \\
\bottomrule
\end{tabular}
\end{table}

在人机协同完成本研究的过程中，作者与人工智能工具的分工如下（详见附录B，完整对话记录见代码仓库）。作者提出了最初关于机载计算与观测如何权衡的问题；在第一版结果出来后要求改用真实卫星数据检验、补充对照与消融实验；在前期研究结果出来后指出用时分配增益过小、方法新意不足，并提出了“两条航路、一张粗图、能否找到让另一条路更好的合理海流反例”这一新问题，决定转向本文的研究。人工智能助手Claude（Anthropic）完成了文献检索、代码、实验与作图，并核对了正文中的数字与引用；Gemini（Google）根据Claude整理的事实清单与提纲撰写了论文文字。
